%%%%%%%%%%%%%%%%%%%%%%%%%%%%%%%%%%%%%%%%%%%%%%%%%%%%%%%%%%%%%%%%%%%%
%% I, the copyright holder of this work, release this work into the
%% public domain. This applies worldwide. In some countries this may
%% not be legally possible; if so: I grant anyone the right to use
%% this work for any purpose, without any conditions, unless such
%% conditions are required by law.
%%%%%%%%%%%%%%%%%%%%%%%%%%%%%%%%%%%%%%%%%%%%%%%%%%%%%%%%%%%%%%%%%%%%

\documentclass{beamer}
\usetheme[faculty=ped,basePath=../fibeamer,logo=../fibeamer/logo/mu/Logo-UC-3.png]{fibeamer}
\graphicspath{{../resources/}}

\usepackage{algorithm}
\usepackage{algpseudocode}

\usetheme[faculty=ped]{fibeamer}
\usepackage[utf8]{inputenc}
\usepackage[
  main=english, %% By using `czech` or `slovak` as the main locale
                %% instead of `english`, you can typeset the
                %% presentation in either Czech or Slovak,
                %% respectively.
  czech, slovak %% The additional keys allow foreign texts to be
]{babel}        %% typeset as follows:
%%
%%   \begin{otherlanguage}{czech}   ... \end{otherlanguage}
%%   \begin{otherlanguage}{slovak}  ... \end{otherlanguage}
%%
%% These macros specify information about the presentation
\title{Variables y expresiones 2} %% that will be typeset on the
\subtitle{Semestre II-2026} %% title page.
\author{Andreina Cota Vidaurre}
%% These additional packages are used within the document:
\usepackage{ragged2e}  % `\justifying` text
\usepackage{booktabs}  % Tables
\usepackage{tabularx}
\usepackage{tikz}      % Diagrams
\usetikzlibrary{calc, shapes, backgrounds}
\usetikzlibrary{positioning, arrows.meta, shadows.blur, shapes.geometric}
\usepackage{amsmath, amssymb}
\usepackage{url}       % `\url`s
\usepackage{listings}  % Code listings
\usepackage{graphicx}
\usepackage{amsmath}
\usepackage[T1]{fontenc}
\usepackage{xcolor}
\usepackage{minted}
% \usepackage{adjustbox}


% Definicion de pseudocodigo en espaniol
\lstdefinelanguage{PseudocodigoES}{
    morekeywords={INICIO,FIN,PARA,HACER,SI,ENTONCES,FIN_SI,FIN_PARA,IMPRIMIR},
    sensitive=true,
    morecomment=[l]{//},
    morestring=[b]"
}

% Definicion de pseudocodigo en ingles
\lstdefinelanguage{PseudocodeEN}{
    morekeywords={BEGIN,END,FOR,DO,IF,THEN,END_IF,END_FOR,PRINT},
    sensitive=true,
    morecomment=[l]{//},
    morestring=[b]"
}

\lstset{
    basicstyle=\ttfamily\small,
    keywordstyle=\bfseries,
    columns=fullflexible,
    keepspaces=true,
    showstringspaces=false,
    frame=single,
    xleftmargin=0em,
    framexleftmargin=0em,
    framesep=2pt,
    gobble=4,
    resetmargins=true,
    aboveskip=0pt,
    belowskip=0pt
}

\lstdefinestyle{pythonstyle}{
    language=Python,
    basicstyle=\ttfamily\small,
    keywordstyle=\bfseries,
    commentstyle=\itshape,
    stringstyle=\ttfamily,
    showstringspaces=false,
    columns=fullflexible,
    keepspaces=true,
    frame=single,
    breaklines=true,
    xleftmargin=0em,
    framexleftmargin=0em,
    framesep=2pt,
    aboveskip=0pt,
    belowskip=0pt,
    gobble=4,
    resetmargins=true,
    emph={print,input,int,float,str,bool,len,range},
    emphstyle=\bfseries
}

\lstset{style=pythonstyle}

% Estilos del diagrama
\tikzstyle{iniciofin} = [
    ellipse,
    minimum width=2.5cm,
    minimum height=1cm,
    text centered,
    draw=black,
    fill=blue!25
]

\tikzstyle{proceso} = [
    rectangle,
    rounded corners,
    minimum width=3cm,
    minimum height=1cm,
    text centered,
    draw=black,
    fill=cyan!20
]

\tikzstyle{decision} = [
    diamond,
    aspect=2,
    text centered,
    draw=black,
    fill=yellow!35
]

\tikzstyle{flecha} = [
    thick,
    ->,
    >=Stealth
]

% \definecolor{green_python}{HTML}{52A736}

% \newcolumntype{Y}{>{\raggedright\arraybackslash}X}

\frenchspacing
\begin{document}
  \shorthandoff{-}
  \frame[c]{\maketitle}

% \begin{darkframes}

    \begin{frame}{Contenido}
        \begin{itemize}
            \item Comparaciones entre strings
            \item \mintinline{python}{print()}
            \item Operadores de asignación \mintinline{python}{+=, -=, *=, /=}
            \item Operadores lógicos \mintinline{python}{and, or, not}
        \end{itemize}
    \end{frame}

    \begin{frame}[fragile,t]{Comparaciones entre strings}
        \begin{itemize}
            \item Se pueden comparar strings usando operadores de comparación
            \item \mintinline{python}{==} verifica si dos strings son iguales
            \item \mintinline{python}{!=} verifica si dos strings son diferentes
            \begin{minted}[xleftmargin=-2cm, frame=leftline, fontsize=\small]{python}
            print("Hola" == "Hola") # True
            print("Hola" != "Hola") # False
            print("Hola" == "Mundo") # False
            print("Hola" != "Mundo") # True
            \end{minted}
            \item Se puede usar \mintinline{python}{==} para verificar si una variable es igual a un string
            \begin{minted}[xleftmargin=-2cm, frame=leftline, fontsize=\small]{python}
            nombre = "Juan"
            print(nombre == "Juan") # True
            \end{minted}
        \end{itemize}
    \end{frame}

    \begin{frame}[fragile,t]{Comparaciones entre strings}
        \begin{itemize}
            \item Python compara los strings carácter por carácter, de izquierda a derecha.
            \item Considera los valores ASCII de los carácteres.
            \item La comparación se detiene en el primer carácter diferente. Se llama orden lexicográfico.
            \item Las letras mayúsculas son consideradas menores que las letras minúsculas.
            \begin{minted}[xleftmargin=-2cm, frame=leftline, fontsize=\small]{python}
            print("Hola" < "Mundo") # True
            print("Hola" > "Mundo") # False
            print("manzana" < "pera")   # True
            print("Zapato" < "manzana")   # ?
            \end{minted}
        \end{itemize}
    \end{frame}

    \begin{frame}[fragile,t]{\mintinline{python}{print()}}
        \begin{itemize}
            \item \mintinline{python}{print()} sin nada dentro simplemente
            imprime un salto de línea.
            \item Puede recibir varios valores separados por comas. 
            Imprime todos los valores en la misma línea, separados por un espacio.
            \begin{minted}[xleftmargin=-2cm, frame=leftline, fontsize=\small]{python}
            print("Hola", "Mundo") # Hola Mundo
            \end{minted}
            \item Se puede concatenar strings con \mintinline{python}{+} dentro del \mintinline{python}{print()}.
            \begin{minted}[xleftmargin=-2cm, frame=leftline, fontsize=\small]{python}
            print("Hola" + "Mundo") # Hola Mundo
            \end{minted}
            \item Diferencia clave: con \mintinline{python}{,} Python agrega espacios automáticamente; 
            con \mintinline{python}{+} hay que agregarlos manualmente y \textbf{todo debe ser texto}.
            \end{minted}
        \end{itemize}
    \end{frame}

    \begin{frame}[fragile,t]{Operadores de asignación}
        \begin{itemize}
            \item Es muy común necesitar modificar el valor de una variable. Puede hacerse
            de forma explícita:
            \begin{minted}[xleftmargin=-2cm, frame=leftline, fontsize=\small]{python}
            contador = 5
            contador = contador + 1
            print(contador) # 6
            \end{minted}
            \item Python ofrece una forma más concisa. Para incrementar el valor se usa \mintinline{python}{+=}:
            \begin{minted}[xleftmargin=-2cm, frame=leftline, fontsize=\small]{python}
            contador = 5
            contador += 1
            print(contador) # 6
            \end{minted}
            \item \mintinline{python}{contador += 1} es equivalente a \mintinline{python}{contador = contador + 1}
        \end{itemize}
    \end{frame}

    \begin{frame}[fragile,t]{Operadores de asignación}
        \begin{itemize}
            \item Se puede decrementar el valor de una variable de la siguiente manera:
            \begin{minted}[xleftmargin=-2cm, frame=leftline, fontsize=\small]{python}
            contador -= 1
            print(contador) # 4
            \end{minted}
            \item Se puede multiplicar el valor de una variable de la siguiente manera:
            \begin{minted}[xleftmargin=-2cm, frame=leftline, fontsize=\small]{python}
            contador *= 2   
            print(contador) # 8
            \end{minted}
        \end{itemize}
    \end{frame}

    \begin{frame}[fragile,t]{Diferencia entre \mintinline{python}{+=} y \mintinline{python}{=+}}
        \begin{itemize}
            \item \mintinline{python}{+=} es un operador de asignación que incrementa el valor de una variable.
            \item \mintinline{python}{=+} incorrecto pero válido en Python. Reasigna, ignorando el valor anterior.
            \item Python interpreta el $+$ como el signo positivo (como escribir $+1$ en vez de 1)
            \item Es equivalente a \mintinline{python}{contador = +1}, es decir, descarta completamente el valor que 
            tenía la variable.
        \end{itemize}
    \end{frame}

    \begin{frame}[fragile,t]{Diferencia entre \mintinline{python}{+=} y \mintinline{python}{=+}}
        \begin{itemize}
            \item Ejemplo de \mintinline{python}{+=}
            \begin{minted}[xleftmargin=-2cm, frame=leftline, fontsize=\small]{python}
            saldo = 100
            saldo += 50      # correcto: saldo pasa a ser 150
            print(saldo)      # 150
            \end{minted}
            \item Ejemplo de \mintinline{python}{=+}
            \begin{minted}[xleftmargin=-2cm, frame=leftline, fontsize=\small]{python}
            saldo = 100
            saldo =+ 50       # error silencioso: ignora el saldo anterior
            print(saldo)       # 50
            \end{minted}
        \end{itemize}
    \end{frame}

    \begin{frame}[fragile,t]{Diferencia entre \mintinline{python}{+=} y \mintinline{python}{=+}}
        \begin{itemize}
            \item Ejemplo de \mintinline{python}{+=} con strings
            \begin{minted}[xleftmargin=-2cm, frame=leftline, fontsize=\small]{python}
            mensaje = "Hola"
            mensaje += " mundo"   # correcto: concatena
            print(mensaje)          # Hola mundo
            \end{minted}
            \item Ejemplo de \mintinline{python}{=+}
            \begin{minted}[xleftmargin=-2cm, frame=leftline, fontsize=\small]{python}
            mensaje = "Hola"
            mensaje =+ " mundo"    # ERROR: no se puede concatenar así
            print(mensaje)
            \end{minted}
            \item En strings se genera un error \mintinline{python}{TypeError}. Intenta 
            aplicar el signo $+$ (unario) a un string, lo cual no es válido.
        \end{itemize}
    \end{frame}

    \begin{frame}[fragile,t]{Tablas de verdad y operadores lógicos}
        \begin{itemize}
            \item La tabla de verdad ayuda a analizar, de forma ordenada, 
            todas las combinaciones posibles entre condiciones y el resultado final de la expresión
            \item Organiza los valores \textbf{Verdadero (V)} y \textbf{Falso (F)} de una o más condiciones.
            \item Muestra el resultado de aplicar operadores lógicos como \mintinline{python}{and, or, not}
            \item Los operadores lógicos permiten construir condiciones más complejas a partir 
            de condiciones simples.
        \end{itemize}
    \end{frame}

    \begin{frame}[fragile,t]{Tablas de verdad: \mintinline{python}{and, or, not}}
        \begin{columns}[T,onlytextwidth]
            \column{0.32\textwidth}
            \centering
            \textbf{\mintinline{python}{and}}\\[0.4em]
            \begin{tabular}{|c|c|c|}
                \hline
                A & B & A and B \\
                \hline
                F & F & F \\
                \hline
                F & V & F \\
                \hline
                V & F & F \\
                \hline
                \textcolor{blue}{V} & \textcolor{blue}{V} & \textcolor{blue}{V} \\
                \hline
            \end{tabular}
    
            \column{0.32\textwidth}
            \centering
            \textbf{\mintinline{python}{or}}\\[0.4em]
            \begin{tabular}{|c|c|c|}
                \hline
                A & B & A or B \\
                \hline
                \textcolor{blue}{F} & \textcolor{blue}{F} & \textcolor{blue}{F} \\
                \hline
                F & V & V \\
                \hline
                V & F & V \\
                \hline
                V & V & V \\
                \hline
            \end{tabular}
    
            \column{0.32\textwidth}
            \centering
            \textbf{\mintinline{python}{not}}\\[0.4em]
            \begin{tabular}{|c|c|}
                \hline
                A & not A \\
                \hline
                F & V \\
                \hline
                V & F \\
                \hline
            \end{tabular}
        \end{columns}
    \end{frame}

    \begin{frame}[fragile,t]{Operador \mintinline{python}{and}}
        \begin{itemize}
            \item La expresión es verdadera si \textbf{ambas condiciones} son verdaderas
            \begin{minted}[xleftmargin=-2cm, frame=leftline, fontsize=\small]{python}
            condicion1 and condicion2
            \end{minted}
            \item Se usa cuando el problema exige que se cumplan \textbf{varios requisitos al mismo tiempo}
            \begin{minted}[xleftmargin=-2cm, frame=leftline, fontsize=\small]{python}
            edad = 20
            tiene_licencia = True

            puede_conducir = edad >= 18 and tiene_licencia
            print(puede_conducir)
            \end{minted}
        \end{itemize}
    \end{frame}

    \begin{frame}[fragile,t]{Operador \mintinline{python}{or}}
        \begin{itemize}
            \item La expresión es verdadera si \textbf{al menos una} condición es verdaderas
            \begin{minted}[xleftmargin=-2cm, frame=leftline, fontsize=\small]{python}
            condicion1 or condicion2
            \end{minted}
            \item Se usa cuando existen \textbf{caminos alternativos} para que algo ocurra
            \begin{minted}[xleftmargin=-2cm, frame=leftline, fontsize=\small]{python}
            tiene_id = False
            tiene_tuc = True
            
            registrado = tiene_id or tiene_tuc
            print(registrado)
            \end{minted}
        \end{itemize}
    \end{frame}

    \begin{frame}[fragile,t]{Operador \mintinline{python}{not}}
        \begin{itemize}
            \item Sirve para \textbf{invertir} el valor lógico de una condición
            \begin{minted}[xleftmargin=-2cm, frame=leftline, fontsize=\small]{python}
            not condicion
            \end{minted}
            \item Ejemplo:
            \begin{minted}[xleftmargin=-2cm, frame=leftline, fontsize=\small]{python}
            valor = False

            valor_invertido = not valor
            print(valor_invertido)
            \end{minted}
        \end{itemize}
    \end{frame}

    \begin{frame}{Mini actividad}
        \begin{itemize}
            \item Un estudiante aprueba una practica si su nota es al menos 4 y su asistencia es al menos 80.
            Qué operador lógico se necesita?
            % \item \mintinline{python}{and}
            \item Un militar puede salir de su puesto si termino su turno o si tiene una autorizacion especial. 
            Qué operador lógico se necesita?
            % \item \mintinline{python}{or}
            \item Un sistema tiene una variable que indica si hay conexión a internet.
            Se quiere obtener el estado contrario de la conexión a internet.
            % \item \mintinline{python}{not}
        \end{itemize}
    \end{frame}

\end{document}
